% ---------------------------------------------------------------------------
% Codebook-regularity paragraph -- drop in at the end of \ref{sub:rtl} (right
% after the fixed-vs-programmable decode discussion) or as a short standalone
% \subsection{Codebook Regularity} before \ref{sec:sw-support}. Cites
% tab:regularity and the Pareto figure gf4_regularity_pareto.png.
% ---------------------------------------------------------------------------
\paragraph{Codebook regularity: why these levels.} The GF4 decode costs
$\approx$1.5\% of the PE (Table~\ref{tab:node-scaling}) precisely because its
Gaussian-quantile magnitudes are scattered-bit constants; the regular E2M1
codebook, whose near-power-of-two levels collapse to a trivial mux, decodes at
$\approx$0.6\%. This invites a question: could GF4's levels be made more regular
--- cheaper to decode, friendlier to shift-based hardware --- without losing
accuracy? We answer it directly by snapping the levels to a Q1.$B$ fixed-point
grid and measuring both axes: WikiText-2 perplexity (OPT-125M, activations
GF4-quantized under a blockwise Hadamard rotation) against decode-ROM area
(synthesized with Yosys against Nangate45). The result (Table~\ref{tab:regularity},
Figure~\ref{fig:gf4-regularity}) has a clear knee. Rounding the levels to Q1.4
(nearest $1/16$) keeps all eight levels distinct, shrinks the decode ROM by
$40\%$ ($15.16\!\rightarrow\!9.04\,\mu\text{m}^2$), and costs only $+0.03$
perplexity --- essentially free. Snapping harder is not: Q1.3 already begins
merging levels ($+0.54$ PPL), and a pure power-of-two codebook collapses three
mid-range levels onto $0.5$, saving more area ($52\%$) but losing $+13$ PPL. The
accuracy-optimal codebook is therefore mildly regularized (Q1.4--Q1.5), not
power-of-two. The pattern replicates across model scale: on OPT-1.3B the same
sweep gives $+0.06$ PPL at Q1.4 and $+13$ at power-of-two, so the knee is a
property of the codebook, not of a particular network.

This tolerance is a direct consequence of GF4's operand form. Because a GF4 value
is $x_i = s_b\,\ell(c_i)$ with per-block scale $s_b = \text{RMS}_b\cdot\alpha$,
the codebook $\ell(\cdot)$ is applied \emph{after} per-block RMS normalization: it
quantizes only the \emph{shape} of the within-block distribution, while the scale
carries the dynamic range. Four bits of relative resolution ($1/16$) are thus
sufficient, and the result is robust to the scale's own precision --- quantizing
the block scale to E4M3, GF4's deployed weight-scale format, moves exact-GF4
perplexity by $+0.05$ and leaves the Q1.4 conclusion unchanged. We retain the
full-precision Gaussian-quantile levels in the reported results, since the decode
is already near-free; the Q1.4 point is available as a drop-in when a design
prioritizes decode area or shift-friendly hardware over the last $0.03$ of
perplexity.

% --- figure float for the Pareto (needs \usepackage{graphicx}) ---
\begin{figure}[t]
  \centering
  \includegraphics[width=0.72\linewidth]{figures/gf4_regularity_pareto.png}
  \caption{Accuracy vs.\ decode-area for regularized GF4 codebooks. Snapping the
  Gaussian-quantile levels to a coarser fixed-point grid shrinks the decode ROM
  (x-axis) but eventually collapses levels and degrades perplexity (y-axis).
  Q1.4 (circled) is the knee: $40\%$ smaller decode for $+0.03$ PPL; power-of-two
  snapping saves more area but costs $+13$ PPL. WikiText-2 / OPT-125M, A-only GF4
  under a blockwise Hadamard.}
  \label{fig:gf4-regularity}
\end{figure}
